\section{结论}
\label{sec:conclusion}

专家拉取、复制和样本重排分别改变通信对象、计算位置和 token 与 expert 的拓扑映射：前者适合权重复用，第二类针对持续性热门 expert，第三类针对路由局部性的跨节点路径。端到端收益取决于主导瓶颈、可隐藏的重配置程度、同步与显存成本及热点持续时间；若动态机制改变 token 接纳或训练轨迹，还须在相同目标质量下用训练总时间补充 step-time。通信调度提供共同执行层优化。后续应在统一模型、拓扑和路由 trace 下分层验证控制变量与真实收敛收益，确定何时移动对象能缩短训练关键路径。

